\chapter{Hall correspondences and primitive kernels}
\label{v4-arith:hall:chapter}

A Hall product counts extensions. Its associativity comes from flags:
the two bracketings forget different steps of the same filtration.
Recovering a geometric operation from the count requires more data.
Automorphisms determine the integration weight, and a vanishing scalar
sum need not come from an acyclic complex.

The first primitive already exhibits both issues. Its coefficients on
a nonsplit extension and a split extension cancel under Green's
coproduct. Placing those coefficients on separate components leaves
nonzero cohomology. A Fourier kernel on the extension line instead
joins the two strata by a Gysin morphism. We construct that kernel,
extend it to every cut of a flag, and determine what is lost when the
flag is forgotten.

Quiver cohomology gives another extension product. Its Euler factors,
cohomological grading and vanishing cycles must be retained before
comparing it with arithmetic fields. An explicit shuffle relation
will also distinguish a valid product from a proposed ordinary
cooperation. The arithmetic comparison remains a problem about these
operations and their maps, beyond equality of numerical invariants.

\section{Extension counts and the tensor crossing}
\label{v4-arith:hall:sec:counts}

Let $\mathcal A$ be an essentially small abelian category with finite
sets $\operatorname{Hom}(M,N)$ and $\operatorname{Ext}^1(M,N)$.
Write $a_M=|\operatorname{Aut}M|$. Define $F^L_{M,N}$ to be the
number of subobjects $N'\subset L$ satisfying
$N'\simeq N$ and $L/N'\simeq M$. Thus the first index names the
quotient and the second names the subobject. Let $H(\mathcal A)$
be the free vector space over a field $k$ on isomorphism classes $u_M$.

\begin{proposition}\label{v4-arith:hall:prop:product}
The formula
\[
 u_Mu_N=\sum_{[L]}F^L_{M,N}u_L
\]
defines an associative multiplication with unit $u_0$.
Let $K_0(\mathcal A)$ be the abelian group generated by object classes,
subject to $[L]=[M]+[N]$ for each short exact sequence.
For any bicharacter
$b:K_0(\mathcal A)\times K_0(\mathcal A)\to k^\times$, the product
$u_M*_b u_N=b([M],[N])u_Mu_N$ is also associative.
Here a bicharacter is multiplicative in each additive argument.
\end{proposition}
\begin{proof}
The finite extension set leaves finitely many possible middle objects.
The finite Hom set bounds the number of embeddings into each one.
The coefficient of $u_L$ in either triple product counts flags
\[
 0\subset L_1\subset L_2\subset L
\]
with successive quotients $P,N,M$. The two bracketings forget
$L_1$ or $L_2$, so their coefficients agree. The twist factors agree
because
\[
 b(M,N)b(M+N,P)=b(N,P)b(M,N+P).
\]
The zero object has precisely the unit subobject and quotient.
\end{proof}

For the coproduct, assume that $\mathcal A$ is linear over
$\mathbb F_q$, hereditary, and has finite-dimensional Hom and Ext
spaces. Hereditary means that $\operatorname{Ext}^i=0$ for $i>1$.
Assume also that every object has finitely many subobjects, as for
finite-dimensional modules over a finite field.
Then
\[
 \langle M,N\rangle
 =\dim_{\mathbb F_q}\operatorname{Hom}(M,N)
  -\dim_{\mathbb F_q}\operatorname{Ext}^1(M,N)
\]
is additive in short exact sequences. Put
$(M,N)=\langle M,N\rangle+\langle N,M\rangle$ and $v=\sqrt q$.
Work over $\mathbb Q(v)$ with the twisted product
$u_M*u_N=v^{\langle M,N\rangle}u_Mu_N$.

\begin{theorem}[Green's compatibility]
\label{v4-arith:hall:thm:green}
The coproduct and counit
\begin{equation}\label{v4-arith:hall:eq:green}
 \Delta u_L=\sum_{[M],[N]}v^{\langle M,N\rangle}
       \frac{a_Ma_N}{a_L}F^L_{M,N}\,u_M\otimes u_N,
 \qquad \varepsilon(u_L)=\delta_{L,0},
\end{equation}
make this Hall algebra a braided bialgebra. For homogeneous classes
$\alpha,\beta$, the crossing is
\[
 c(x_\alpha\otimes y_\beta)
   =v^{(\alpha,\beta)}y_\beta\otimes x_\alpha.
\]
In particular, multiplicativity of $\Delta$ uses
\begin{equation}\label{v4-arith:hall:eq:crossing}
 (x_\alpha\otimes y_\beta)(x_\gamma\otimes y_\delta)
 =v^{(\beta,\gamma)}(x_\alpha*x_\gamma)
                     \otimes(y_\beta*y_\delta).
\end{equation}
All sums in these formulas are finite.
\end{theorem}
\begin{proof}
The automorphism factors in an iterated coproduct cancel, leaving
the same three-step flags in either order. This proves coassociativity.

For multiplicativity, Green's formula counts pairs $N,N'\subset L$
through their four subquotients
\[
 L/(N+N'),\quad N/(N\cap N'),\quad
 N'/(N\cap N'),\quad N\cap N'.
\]
If the first and last are $A,D$, the groupoid of extension-square
fillings contributes
\[
 \frac{|\operatorname{Ext}^1(A,D)|}
      {|\operatorname{Hom}(A,D)|}
 =q^{-\langle A,D\rangle}.
\]
Heredity removes the degree-two obstruction to a filling.


Substituting these weights into \eqref{v4-arith:hall:eq:green} leaves
the crossing \eqref{v4-arith:hall:eq:crossing}, by Euler additivity.
This is Green's extension-counting theorem with its automorphism
normalization. See Ringel, \emph{Green's theorem on Hall algebras},
Part~I, \S\S1 and~8. The zero-object terms prove both counit identities.
\end{proof}

For finite vector spaces let $s=u_{\mathbb F_q}$. Counting lines gives
\[
 s*s=v(q+1)u_{\mathbb F_q^2},\qquad
 [s\otimes s]\Delta u_{\mathbb F_q^2}=v^{-1}.
\]

Here $[s\otimes s]$ denotes the coefficient of that basis vector.
Thus $[s\otimes s]\Delta(s*s)=q+1$. Expanding $\Delta(s)^2$
gives $1+q$, since the crossed term has weight $v^2=q$.
An ordinary flip would give $2$. Even this smallest example fixes
the tensor multiplication used by the theorem.

A hereditary category can have infinitely many subobjects, as for
coherent sheaves on a curve. Its coproduct requires a specified tensor
completion and locally finite products. The finite extension-square
identity persists. It does not by itself define multiplication on
arbitrary coefficientwise infinite sums.

The hereditary hypothesis also has content. For
$R=\mathbb F_q[\epsilon]/(\epsilon^2)$ and $S=R/(\epsilon)$,
the resolution with every differential multiplication by $\epsilon$
gives $\operatorname{Ext}^n_R(S,S)=\mathbb F_q$ for every $n\geq0$.
The alternating sums of Ext dimensions oscillate. The Hall product
still exists, but this infinite Ext sequence does not define a finite
Euler twist.

\section{The local primitive and its Green norm}
\label{v4-arith:hall:sec:primitive}

Fix a prime power $Q$ and the category of finite-dimensional
nilpotent $\mathbb F_Q[t]$-modules. A partition
$\lambda=(\lambda_1,\ldots,\lambda_l)$ names the module
$M_\lambda=\bigoplus_i\mathbb F_Q[t]/(t^{\lambda_i})$.
Write $u_\lambda=u_{M_\lambda}$ and $|\lambda|=\sum_i\lambda_i$.
The length $\ell(\lambda)$ is its number of Jordan blocks.
The Euler form is zero, so Green's multiplication and coproduct
use the ordinary tensor product here.

Work over $\mathbb C$ with the positive Hermitian Green form
\[
 (u_\lambda,u_\mu)_G
 =\frac{\delta_{\lambda\mu}}{a_{M_\lambda}},
\]
conjugate-linear in the first argument. The coproduct is adjoint
to multiplication. For $r\geq0$, put
$e_r=Q^{r(r-1)/2}u_{(1^r)}$, with $e_0=1$, and set
$E(z)=\sum_{r\geq0}e_rz^r$. Subspace counting gives
\[
 \Delta u_{(1^r)}=\sum_{i+j=r}Q^{-ij}u_{(1^i)}\otimes u_{(1^j)},
 \qquad \Delta E(z)=E(z)\otimes E(z).
\]
These are formal power series, with finite-dimensional pieces in
each module dimension.

\begin{theorem}[Classical Hall primitives]
\label{v4-arith:hall:thm:primitives}
The algebra is polynomial on the $e_r$, $r>0$.
Define $p_r$ by
\begin{equation}\label{v4-arith:hall:eq:newton}
 \log E(z)=\sum_{r\geq1}\frac{(-1)^{r-1}}r p_rz^r.
\end{equation}
Then
\begin{equation}\label{v4-arith:hall:eq:partitions}
 p_r=\sum_{|\lambda|=r}
       \left(\prod_{j=1}^{\ell(\lambda)-1}(1-Q^j)\right)u_\lambda.
\end{equation}
Each $p_r$ is primitive, and the homogeneous primitives of dimension
$r$ form its one-dimensional span. Their Green norms are
\begin{equation}\label{v4-arith:hall:eq:norm}
 (p_r,p_s)_G=\delta_{rs}\frac{r}{Q^r-1}.
\end{equation}
\end{theorem}
\begin{proof}
The polynomial presentation is the classical Hall theorem.
Its triangular construction uses products of semisimple classes
whose dimensions are the successive kernel increments of $t$.
The leading term has the kernel filtration of $M_\lambda$.
Every other term has larger kernel dimensions. In each fixed
dimension the resulting triangular matrix is invertible.
Commutativity follows by dualizing exact sequences, since every
nilpotent Jordan module is isomorphic to its linear dual.

The coefficient formula \eqref{v4-arith:hall:eq:partitions} is the
classical Newton identity in this basis. In the present normalization,
it is Lemma~3.5 of Shimoji--Yanagida,
\emph{A study of symmetric functions via derived Hall algebra},
arXiv:1812.06033. Its coefficients are fixed by the vertical-strip
Hall numbers, rather than by rescaling a primitive after construction.
Taking the logarithm of the group-like identity for $E$ proves
primitivity. The Newton relations identify the polynomial generators
$e_r$ and $p_r$ triangularly.

In a characteristic-zero polynomial algebra with primitive
generators, a primitive polynomial $f$ satisfies $f(X+Y)=f(X)+f(Y)$.
Differentiating in each $Y$ variable at zero makes every partial
derivative of $f$ constant. Hence $f$ is linear in the $p_r$.
The dimension grading gives the asserted homogeneous span.

Finally, the order of the general linear group gives
\[
 (E(u),E(v))_G
 =\sum_{r\geq0}
   \frac{Q^{r(r-1)/2}(\bar uv)^r}{\prod_{j=1}^r(Q^j-1)}
 =\prod_{j\geq1}(1-Q^{-j}\bar uv)^{-1}.
\]
The product identity follows coefficientwise from
$(1-t)P(Qt)=P(t)$ and $P(0)=1$.
Primitivity and adjunction make distinct primitive dimensions
orthogonal and give factorial pairings for powers. Taking logarithms
of this identity and using \eqref{v4-arith:hall:eq:newton} yields
\eqref{v4-arith:hall:eq:norm}.
\end{proof}

In dimension two, writing $S=\mathbb F_Q[t]/(t)$ and
$J_2=\mathbb F_Q[t]/(t^2)$ gives
\[
 u_S^2=u_{J_2}+(Q+1)u_{S\oplus S},\qquad
 p_2=u_{J_2}-(Q-1)u_{S\oplus S}.
\]
The middle Green coefficients of $u_{J_2}$ and $u_{S\oplus S}$
are $(Q-1)/Q$ and $1/Q$. They cancel for $p_2$.
The cancellation uses automorphism orders and a signed linear
combination of classes. It does not yet specify a complex or a
homotopy on an extension space.

\section{Flag convolution and a nonacyclic primitive class}
\label{v4-arith:hall:sec:groupoids}

Let $\mathcal M_n$ be a finite skeleton of the groupoid of
$n$-dimensional nilpotent $\mathbb F_Q[t]$-modules.
For a characteristic-zero field $k$, consider
\[
 \mathscr D_n=D^b\operatorname{Fun}
       (\mathcal M_n,\operatorname{Vect}^{\mathrm{fd}}_k)
 \simeq\prod_{|\lambda|=n}
       D^b\operatorname{Rep}^{\mathrm{fd}}_k(\operatorname{Aut}M_\lambda).
\]
A functor on a groupoid assigns a vector space to each object and
an invertible map to each isomorphism. Its restriction to one
component is a representation of that object's automorphism group.
Let $\mathcal E_{m,n}$ be the groupoid of short exact sequences
with quotient dimension $m$ and subobject dimension $n$.
Write $p$ for its endpoint map and $q$ for its middle-object map.

Pullback followed by left Kan extension defines
\[
 \mathscr F\diamond\mathscr G
 =q_!p^*(\mathscr F\boxtimes\mathscr G).
\]
Left Kan extension sums over the relevant finite groupoid fiber,
taking coinvariants under its stabilizers. Characteristic zero makes
these finite-group coinvariants exact. The two iterated convolution
functors are push--pull along the same three-step flag groupoid.
Base change gives associativity. Full flags give its higher coherences.

Let $\delta_M$ be the trivial representation on the component of
$M$, extended by zero. At $L$, the fiber of
$\delta_M\diamond\delta_N$ is the permutation representation on
subobjects counted by $F^L_{M,N}$. Taking dimensions recovers the
Hall product. For the reverse functor $p_!q^*$, coinvariants of the trivial
representation of a finite group have dimension one. Groupoid
cardinality instead assigns the reciprocal group order.

\begin{proposition}\label{v4-arith:hall:prop:virtual}
The zero-differential complex
\[
 \mathscr P_2=\delta_{J_2}
       \oplus\delta_{S\oplus S}^{\oplus(Q-1)}[1]
\]
has Euler class $p_2$. After applying $p_!q^*$ and forgetting
endpoint actions, its $(S,S)$ component has
\[
 H^0=k^{Q-1},\qquad H^{-1}=k^{Q-1}.
\]
It is not acyclic.
\end{proposition}
\begin{proof}
There are $Q$ extension classes of $S$ by $S$.
Each has stabilizer $\operatorname{Hom}(S,S)$ of order $Q$.
The $Q-1$ nonzero classes have middle module $J_2$. The remaining
class is split. Ordinary Kan extension therefore gives dimensions
$Q-1$ and $1$. It preserves the zero differential, so the shifted
split summands give exactly the two displayed cohomology groups.


Green's scalar coefficients instead cancel as
$(Q-1)/Q-(Q-1)/Q=0$.
\end{proof}

The two source summands lie on different connected components of
$\mathcal M_2$. This category has no morphism between them that
could provide a cancelling differential. A geometric extension
space can connect the split point to nonsplit extensions.
The next construction retains that connection and its integration
weight at the level of complexes.

\section{A Fourier kernel on the extension space}
\label{v4-arith:hall:sec:fourier}

Fix $\mathbb F_q$, a prime $\ell\nmid q$, and
$K=\overline{\mathbb Q}_\ell$.
We use bounded constructible complexes with Weil structure.
Thus Frobenius acts on their stalk cohomology, and its alternating
trace is defined. Our Tate convention gives $K(1)$ trace $q^{-1}$.
A cohomological shift $[1]$ changes the trace sign.
The notation $Rf_!$ denotes the derived direct image with proper
supports, and $Rf_*$ the ordinary derived direct image.

Choose a nontrivial additive character
$\psi:\mathbb F_q\to K^\times$. Its Artin--Schreier sheaf
$\mathcal L_\psi$ is the rank-one summand of the finite additive
cover selected by $\psi$. For a $d$-dimensional vector space
$V$, $d>0$, let $p_V:V\times(V^\vee\setminus0)\to V$ be projection.
The pairing defines
\begin{equation}\label{v4-arith:hall:eq:fourier-kernel}
 \mathcal P_V
 =Rp_{V!}\mathcal L_{\psi\langle v,\xi\rangle}[1].
\end{equation}
The shift is $[1]$ in every dimension. We use unshifted Fourier
integration before inserting this shift.

\begin{theorem}\label{v4-arith:hall:thm:fourier-primitive}
Let $i:\{0\}\hookrightarrow V$ and
$j:V\setminus0\hookrightarrow V$. There is a triangle
\begin{equation}\label{v4-arith:hall:eq:gysin}
 i_*K(-d)[-2d]\xrightarrow{\mathrm{Gys}}K_V
 \longrightarrow\mathcal P_V\longrightarrow.
\end{equation}
The first arrow is the Gysin map normalized by compact trace,
and $\mathcal P_V\simeq Rj_*K_{V\setminus0}$.
Its nonzero stalk cohomology is
\[
 H^0(\mathcal P_V)_v=K\quad(v\in V),\qquad
 H^{2d-1}(\mathcal P_V)_0=K(-d).
\]
Its trace is $1$ away from zero and $1-q^d$ at zero, while
$R\Gamma_c(V,\mathcal P_V)=0$.
\end{theorem}
\begin{proof}
Fourier integration over all of $V^\vee$ gives
$i_*K(-d)[-2d]$. Indeed, a nonzero $v$ supplies a nonconstant
linear coordinate whose Artin--Schreier sheaf has zero compact
cohomology. At $v=0$, the coefficient is constant on $\mathbb A^d$.
Localization at the zero covector gives
\eqref{v4-arith:hall:eq:gysin}, with the first arrow initially
defined by this integral.

Integrate in the other order. Each $\xi\ne0$ again supplies a
nonconstant Artin--Schreier coordinate, so the entire compact
cohomology vanishes. The first arrow consequently induces the
trace isomorphism on compact cohomology.


Purity identifies it
with the normalized Gysin generator. The localization triangle
$i_*i^!K_V\to K_V\to Rj_*j^*K_V$ identifies its cone.
At zero, the two terms have degrees $2d$ and $0$. Their cone
has the two stated degrees. The odd degree $2d-1$ contributes
$-q^d$ to the trace.
\end{proof}

The vanishing used in this proof is a vanishing of complexes.
For the Artin--Schreier cover of $\mathbb A^1$, its deck translations
act trivially on the sole compact cohomology of the covering affine
line. Every nontrivial character summand therefore vanishes.
Splitting off a nonconstant coordinate proves the higher-dimensional
case. This also fixes the normalization independently of a character
sum. Compare Laumon's Fourier formalism, \S\S1.1.3 and~1.2.1.

For the extension line of $S$ by $S$, the operator is
\[
 \begin{pmatrix}0&v\\0&0\end{pmatrix},\qquad v\in\mathbb A^1.
\]
Its middle module is $J_2$ when $v\ne0$ and $S\oplus S$ when
$v=0$. Thus the trace of $\mathcal P_{\mathbb A^1}$ gives the
coefficients of $p_2$, and integration over the extension line is
zero as a complex. The Gysin arrow supplies the information absent
from the disconnected complex of Proposition~\ref{v4-arith:hall:prop:virtual}.

The construction works for a vector bundle $E\to B$ by using
$E\times_BE^\vee$ and removing the zero section of the dual.
It commutes with base change. Bundle automorphisms act
contragrediently on the dual and preserve the pairing. The kernel
therefore carries equivariance for the endpoint gauge action.
This does not yet make it an object on the stack of middle modules.

\begin{proposition}\label{v4-arith:hall:prop:stabilizer}
For the split vector group $U=\mathbb G_a^h$,
\[
 R\Gamma_c(BU,K)=K(h)[2h].
\]
Its trace is $q^{-h}$. On a product of an extension space with
$BU$, proper-support integration includes this tensor factor.
\end{proposition}
\begin{proof}
The smooth atlas $\operatorname{Spec}\mathbb F_q\to BU$ has
\v Cech nerve $U^n$. Each term has ordinary cohomology $K$.
The normalized complex of these constants has only $K$ in degree zero.
The smooth stack has dimension $-h$ and dualizing complex
$K(-h)[-2h]$. Verdier duality therefore gives the displayed compact
cohomology. The product assertion follows from proper-support
K\"unneth.
\end{proof}

For extensions of $S$ by $S$, automorphisms fixing the endpoints
form $\mathbb G_a$. The nonsplit and split trace contributions are
$(q-1)/q$ and $(1-q)/q$.
The attaching map makes their compact pushforward zero before and
after tensoring with $K(1)[2]$. This geometric factor recovers the
reciprocal stabilizer weight that ordinary finite-group coinvariants
do not supply.

\section{Every cut of a flag}
\label{v4-arith:hall:sec:flags}

For a vector space $E$, write the unshifted Fourier transform as
\[
 \mathsf F_E(\mathcal K)=R\operatorname{pr}_{E^\vee!}
   (\operatorname{pr}_E^*\mathcal K
       \otimes\mathcal L_{\psi\langle x,\xi\rangle}).
\]
Integrating twice gives pullback by $-1$ followed by
$(-\dim E)[-2\dim E]$. The rank-normalized transform is
$\mathsf F_E[\dim E]$.

Let $0=F_0\subset F_1\subset\cdots\subset F_s=V$ be a flag
with positive graded dimensions $n_i$. Define
\[
 \mathfrak g=\operatorname{End}(V),\quad
 \mathfrak p=\{x:xF_i\subset F_i\},\quad
 U=\{x:xF_i\subset F_{i-1}\},\quad
 \mathfrak l=\bigoplus_i\operatorname{End}(F_i/F_{i-1}).
\]
Let $i:\mathfrak p\hookrightarrow\mathfrak g$ and
$q_0:\mathfrak p\to\mathfrak l$ be the inclusion and quotient.
Put $r=\dim U=\sum_{i<j}n_in_j$.
Trace pairings identify the full matrix spaces with their duals.

\begin{proposition}\label{v4-arith:hall:prop:parabolic-fourier}
For every constructible complex $\mathcal K$ on $\mathfrak l$,
\[
 \mathsf F_{\mathfrak g}(i_*q_0^*\mathcal K)
 \simeq i_*q_0^*\mathsf F_{\mathfrak l}(\mathcal K)(-r)[-2r].
\]
For rank-normalized transforms, only the Tate twist remains.
The comparison commutes with base change of flags, proper forgetting
of flags, and finite flag refinements.
\end{proposition}
\begin{proof}
Choose an adapted frame. Integrating the strictly upper blocks
forces the corresponding lower blocks of the dual matrix to vanish.
This says $U^\perp=\mathfrak p$ for the trace pairing.
The affine integral on this locus contributes $K(-r)[-2r]$.
The remaining phase is the trace pairing on $\mathfrak l$.
Since $\dim\mathfrak g=\dim\mathfrak l+2r$, rank normalization
removes exactly that shift.

The pairings and subbundles are intrinsic, so the calculation glues
over the flag scheme. Its incidence projection to $\mathfrak g$
is proper. Proper-support base change commutes the projection with
Fourier integration. Every refinement integrates each strictly upper
block once, with the same phase. Fubini and the coherent composition
of the six operations identify all parenthesizations.
\end{proof}

Full matrix support matters here. The Fourier transform of a
skyscraper at zero is constant on the dual matrix space. It need
not retain nilpotent support.

For $1\leq p<s$, set
\[
 U_p=\{x:x(V)\subset F_p,\ x(F_p)=0\}\subset U,
 \qquad H_p=U_p^\perp\subset U^\vee.
\]
In an adapted splitting, $U_p$ consists of blocks $i\leq p<j$.


Define
\begin{equation}\label{v4-arith:hall:eq:all-cuts}
 \Omega=U^\vee\setminus\bigcup_{p=1}^{s-1}H_p,
 \qquad
 \mathcal P_F=\mathsf F_{U^\vee}(j_!K_\Omega)[s-1],
\end{equation}
where $j:\Omega\hookrightarrow U^\vee$.
A covector belongs to $\Omega$ exactly when its restriction to
each whole cross-cut subspace is nonzero.
For $s=1$, set $\mathcal P_F=K$ on $U=0$.
This is an operation unit. A vacuum at a retained position is
additional data.

\begin{theorem}\label{v4-arith:hall:thm:cut-vanishing}
Every cut quotient $\rho_p:U\to U/U_p$ satisfies
$R\rho_{p!}\mathcal P_F=0$.
The kernel and these comparisons are equivariant for the
flag-preserving group. Successive quotient integrations have
compatible Fubini comparisons.
\end{theorem}
\begin{proof}
Integrate first along $U_p$. A covector nonzero on $U_p$ gives
zero compact cohomology. A covector zero on $U_p$ belongs to
$H_p$, where $j_!K_\Omega$ restricts to zero.
Proper-support base change proves the vanishing on the quotient.
The flag group preserves each $U_p$ and its annihilator, hence
preserves the kernel. Iterated integration integrates over the
same sum of cut subspaces, so its comparisons agree.
\end{proof}

The full flag is visible even in its scalar trace. For
$J\subset\{1,\ldots,s-1\}$, put
$S_J=\sum_{p\in J}U_p$, $r_J=\dim S_J$ and $D=\dim U$.
Inclusion--exclusion and character orthogonality give
\begin{equation}\label{v4-arith:hall:eq:flag-trace}
 k_F(x)=(-1)^{s-1}\sum_J(-1)^{|J|}q^{D-r_J}1_{S_J}(x).
\end{equation}

For three one-dimensional pieces, with coordinates
$(a,b,c)=(x_{12},x_{13},x_{23})$, this becomes
\[
 k_3(x)=q^3 1_{x=0}-q1_{c=0}-q1_{a=0}+1.
\]
A nonzero outer entry $b$, with $a=c=0$, gives $1-2q$.
A single nonzero adjacent entry gives $1-q$.
Both operators have Jordan type $(2,1)$, so the trace remembers
the flag beyond the middle module.

For four one-dimensional pieces let $\delta_{ij}=1_{x_{ij}=0}$.
Formula~\eqref{v4-arith:hall:eq:flag-trace} gives
\[
 \begin{split}
 k_4={}&1-q(\delta_{12}+\delta_{23}+\delta_{34})
        +q^2\delta_{12}\delta_{34}\\
 &+q^3(\delta_{23}\delta_{24}\delta_{34}
          +\delta_{12}\delta_{13}\delta_{23})-q^6 1_{x=0}.
 \end{split}
\]
Meeting every cut is weaker than connecting the graph of labels.
The two edges $13$ and $24$ meet every consecutive cut but
form a disconnected graph.

\subsection{Comparison with the scalar Hall primitive}

For a complete flag in dimension $n$, define
\[
 p_n(x)=\prod_{j=1}^{\ell(\lambda(x))-1}(1-q^j),\qquad x\in U,
\]
where $\lambda(x)$ is the Jordan partition.
These are the coefficients in \eqref{v4-arith:hall:eq:partitions}.
For functions on the dual space, use
$\mathsf Ff(x)=\sum_\xi f(\xi)\psi(\langle x,\xi\rangle)$.
Its inverse has opposite phase and factor $q^{-D}$.

\begin{proposition}\label{v4-arith:hall:prop:primitive-trace}
For each proper cut and every $x\in U$,
\[
 \sum_{u\in U_p(\mathbb F_q)}p_n(x+u)=0.
\]
Hence $\mathsf F^{-1}p_n$ is supported on $\Omega$.
For $n=3$, let $i_{13}:L_{13}\hookrightarrow U$ be the line
$a=c=0$, and write $\mathcal P=\mathcal P_{\mathbb A^1}$.
Then
\begin{equation}\label{v4-arith:hall:eq:outer-correction}
 p_3=\operatorname{tr}
       (\mathcal P_F\oplus i_{13*}\mathcal P(-1)[-2]).
\end{equation}
Both summands have zero proper-support image under either cut.
In dual coordinates $(\alpha,\beta,\gamma)$,
\[
 p_3=\mathsf F(1_\Omega-q^{-1}1_{\{\beta\ne0\}}).
\]
\end{proposition}
\begin{proof}
Fix the diagonal blocks $A,C$ at a cut. Changing a splitting
changes its cross block $B$ by $AZ-ZC$.
The cokernel is $\operatorname{Ext}^1_{\mathbb F_q[t]}(C,A)$,
so every extension class has the same number of cross-block
representatives. The Green coefficient is the sum of the primitive's
middle-object coefficients over these classes, divided by
$|\operatorname{Hom}_{\mathbb F_q[t]}(C,A)|$.
Primitivity makes it zero. Fourier transformation identifies this
vanishing with zero restriction to $H_p$.

For three blocks, $ac\ne0$ gives type $(3)$, a nonzero operator
with $ac=0$ gives type $(2,1)$, and zero gives type $(1,1,1)$.
Their primitive coefficients are $1$, $1-q$ and $(1-q)(1-q^2)$.
Subtracting from $k_3$ gives
\[
 k_3-p_3=-q1_{L_{13}}+q^2 1_{\{0\}}
         =-q\operatorname{tr}(i_{13*}\mathcal P).
\]

The twist $(-1)$ multiplies trace by $q$ and $[-2]$ preserves
its sign. The invariant line lies in both cut subspaces.
Its cut pushforward has fiber $R\Gamma_c(\mathbb A^1,\mathcal P)=0$.
Finally, summing first over $\alpha,\gamma$ gives
\[
 \mathsf F1_{\{\beta\ne0\}}
 =q^2 1_{L_{13}}(q1_{b=0}-1),
\]
which proves the last formula.
\end{proof}

The direct sum in \eqref{v4-arith:hall:eq:outer-correction} realizes
the scalar primitive and its cut vanishing. It supplies no
induction-compatible comparison of Hall products.
The trace identity also supplies no distinguished map whose cone
is the outer-line summand.

\section{The attaching maps and their contractions}
\label{v4-arith:hall:sec:gysin-cube}

The flag kernel also has a construction over the integral flag
scheme. Use constructible rational Beilinson motives, with their
oriented six operations and purity for regular closed immersions.
These operations and realizations are the setting of
Cisinski--D\'eglise, \emph{Triangulated categories of mixed motives},
Theorems~14.4.1 and~15.2.1 and Chapter~17.
No integral Artin--Schreier sheaf is being assumed.

Over the flag scheme, $U_p$ and $S_J$ are vector subbundles of
$U$: their adapted-frame descriptions have constant rank and glue.
Let $i_J:S_J\hookrightarrow U$ and $h_J=D-r_J$. Define
\begin{equation}\label{v4-arith:hall:eq:cube-vertices}
 B_J=i_{J*}\mathbf1_{S_J}(-h_J)[-2h_J].
\end{equation}
Purity identifies $B_J$ with $i_{J*}i_J^!\mathbf1_U$.

For $p\notin J$, the counit of the nested closed immersion gives
a face $B_J\to B_{J\cup\{p\}}$.
If its codimension is $c$, this is the Gysin map
$\mathbf1_{S_J}(-c)[-2c]\to\mathbf1_{S_{J\cup\{p\}}}$,
with closed direct image understood.
Composition of the counits gives a coherent commuting cube.

Place $B_J$ in external degree $|J|$, with total graded term
$B_J[-|J|]$. The face in direction $p$ has sign
\[
 \epsilon(J,p)=(-1)^{\#\{j\in J:j<p\}}.
\]
The total differential is
$d_{\mathrm{Tot}}=(-1)^{|J|}d_{\mathrm{int}}+d_{\mathrm{cube}}$.
Use a compatible complex model of the finite cube and put
$\mathcal M_F=\operatorname{Tot}(B_\bullet)[s-1]$.
The Gysin morphisms remain part of this totalization.

\begin{proposition}\label{v4-arith:hall:prop:realizations}
The $\ell$-adic realization of $\mathcal M_F$ on a finite-field
fiber is $\mathcal P_F$.
The complex realization is the same total Gysin cube, retaining
the Tate markings. On the fiber over a fixed complex flag, use
the right differential-module convention
$\operatorname{DR}(\omega_S)=\mathbb C_S[\dim S]$. Its vertex is
\begin{equation}\label{v4-arith:hall:eq:normal-module}
 (i_J)_+\omega_{S_J}[-r_J-2h_J],
 \qquad\text{with Tate marking }(-h_J).
\end{equation}
\end{proposition}
\begin{proof}
On $U^\vee$, the tensor product of localization fibers
$[K_{U^\vee}\to K_{H_p}]$ is $j_!K_\Omega$.
Its $J$-vertex is $K_{H_J}$ in external degree $|J|$, where
$H_J=\bigcap_{p\in J}H_p=S_J^\perp$.
Fourier integration of this vertex gives $B_J$: the integral
vanishes off $S_J$ and contributes $(-h_J)[-2h_J]$ on it.
Each restriction face becomes the Gysin face normalized by compact
trace. Thus Fourier integration identifies the entire cubes,
including their total differentials and final shift.

Realization preserves the six operations and fundamental classes.
Applying $\operatorname{DR}$ to
\eqref{v4-arith:hall:eq:normal-module} gives
$i_{J*}\mathbb C_{S_J}[-2h_J]$, as required.
For the absolute differential module over a smooth complex base
$B$ of flags, the shift is instead $[-\dim B-r_J-2h_J]$,
because $\dim S_J=\dim B+r_J$.
For a normal coordinate $t$, the sequence
\[
 0\longrightarrow\omega\longrightarrow\omega(*t)
   \longrightarrow i_+\omega_{t=0}\longrightarrow0
\]
represents the codimension-one Gysin map after these shifts.
Tensoring its normal-coordinate models gives higher codimension.
\end{proof}

The supported differential module retains every finite normal
derivative. If $\delta$ is its normal generator, then
\[
 (\delta\partial_t^a)t=a\delta\partial_t^{a-1}.
\]
These are finite sums of derivatives, with no product completion.
The comparison identifies the constructed cube in two realizations.
It does not identify their whole coefficient categories.

\begin{theorem}\label{v4-arith:hall:thm:contractions}
After quotient integration in a nonempty set $P$ of cut directions,
the total Gysin cube admits homotopies $h_p$, $p\in P$, with
\[
 dh_p+h_pd=1,\qquad h_p^2=0,\qquad
 h_ph_q+h_qh_p=0\quad(p\ne q).
\]
They are compatible with the trace-normalized comparisons for
further cut integrations.
\end{theorem}
\begin{proof}
Set $U_P=\sum_{p\in P}U_p$ and let $\rho_P:U\to U/U_P$.
The pushed vertex is supported on $\overline S_J=\rho_P(S_J)$.
Its twist and shift are determined by
\[
 h_J+\dim(S_J\cap U_P)=D-\dim\overline S_J.
\]
Both this integer and the support depend only on $J\setminus P$.

For $p\in P$, affine trace identifies the pushed $p$-face with
the identity. The remaining faces are Gysin maps of the image
subbundles, with common factor $(-\dim U_P)[-2\dim U_P]$.
The pushed cube is therefore equivalent to the remaining cube
tensored with $[K\xrightarrow{1}K]$ in every direction of $P$.

Write $e_J$ for the external label of a vertex in this tensor model.
Before the final shift define
\[
 h_p(e_J)=
 \begin{cases}
 \epsilon(J\setminus\{p\},p)e_{J\setminus\{p\}},&p\in J,\\
 0,&p\notin J.
 \end{cases}
\]
Its coefficient map is the inverse of the identity comparison.

The $p$-face contributes the identity to $dh_p+h_pd$.
Every other face cancels with the reversed order, and the internal
differential cancels because label removal changes $|J|$ by one.
Removing the same label twice gives zero. Reversing two distinct
removals changes the sign. This proves all three identities.

The final shift multiplies both $d$ and $h_p$ by $(-1)^{s-1}$.
Composing affine traces gives the trace for the composite quotient.
The common orientation therefore makes these homotopies compatible
with further integrations.
\end{proof}

\subsection{The omitted outer covector locus}

For three one-dimensional pieces, the two binary groupings use
the dual opens
\[
 O_L=\{\alpha\ne0,\ (\beta,\gamma)\ne(0,0)\},\qquad
 O_R=\{\gamma\ne0,\ (\alpha,\beta)\ne(0,0)\}.
\]
Their union $O$ misses the nonempty locus
$Z=\{\alpha=\gamma=0,\ \beta\ne0\}$ in $\Omega$.
Extending by zero to $U^\vee$, the two triangles are
\[
 \begin{gathered}
 K_{O_L\cap O_R,!}\longrightarrow K_{O_L,!}\oplus K_{O_R,!}
       \longrightarrow K_{O,!}\longrightarrow,\\
 K_{O,!}\longrightarrow K_{\Omega,!}
       \longrightarrow K_{Z,!}\longrightarrow.
 \end{gathered}
\]

Fourier transform with the common shift $[2]$ preserves them.
The last term has trace $q1_{b=0}-1$, so it is nonzero.
These are extension coordinates. A position-space $13|2$ operation
requires a separate map preserving the corresponding supports.

\section{Descent after forgetting the flag}
\label{v4-arith:hall:sec:descent}

Assume the finite field has characteristic different from two.
Let $\mathcal N_2\subset\mathfrak{sl}_2$ be the nilpotent cone,
with punctured inclusion $j_N$.
The extension line maps into it by
$f(v)=\left(\begin{smallmatrix}0&v\\0&0\end{smallmatrix}\right)$.

\begin{theorem}\label{v4-arith:hall:thm:marked-no-descent}
There is no constructible complex $\mathcal F$ on $\mathcal N_2$
with $j_N^*\mathcal F\simeq K$ and
$f^*\mathcal F\simeq\mathcal P_{\mathbb A^1}$, whose restriction
to $\mathbb G_m$ agrees with that constant marking.
\end{theorem}
\begin{proof}
The cone is $\mathbb A^2/\{\pm1\}$, and the quotient is finite
\'etale away from zero. Characteristic-zero coefficients make
invariants exact. The scalar involution preserves the orientation
class, so purity gives stalk cohomology $K$ in degree zero and
$K(-2)$ in degree three for $Rj_{N*}K$ at zero.
In particular, its degree-one stalk vanishes.

An $\mathcal F$ as stated would give, by adjunction and base change,
\[
 \mathcal P_{\mathbb A^1}\longrightarrow f^*Rj_{N*}K
       \longrightarrow Rj_*K_{\mathbb G_m}
       \simeq\mathcal P_{\mathbb A^1}.
\]
The composite is the identity on $\mathbb G_m$.
Since $Rj_*$ is fully faithful, it is the identity everywhere.
Its degree-one stalk would factor the identity of $K(-1)$ through
zero, a contradiction.
\end{proof}

There is a further obstruction without an open marking.
Work over $\overline{\mathbb F}_q$ and let $G=\mathrm{GL}_n$.
For a parabolic $P$ with Levi quotient $L$, use
\[
 \mathcal N_L\xleftarrow{p}\mathcal N_P\xrightarrow{i}\mathcal N_n,
 \qquad \mathcal N_P=\mathcal N_n\cap\operatorname{Lie}P.
\]
Parabolic restriction $\mathsf R_L^G=Rp_!i^*$ first forgets
equivariance from $G$ to $P$, then identifies $P$- and
$L$-equivariance on $\mathcal N_L$. Use its perverse-exact
normalization, and fix a maximal torus $T\subset B\subset G$.

\begin{proposition}\label{v4-arith:hall:prop:unmarked-detection}
If $\mathcal F\in D^b_G(\mathcal N_n,K)$ and
$\mathsf R_T^G\mathcal F=0$, then $\mathcal F=0$.
For $n\geq2$, no nonzero object of this category has all proper
parabolic restrictions zero.
\end{proposition}
\begin{proof}
Parabolic restriction is perverse-exact and has parabolic induction
as right adjoint. See Achar--Mautner, arXiv:1207.7044,
Remark~2.1 and Proposition~3.1.
Every simple $G$-equivariant perverse sheaf is a summand of the
Springer sheaf $\mathsf I_T^G K$.
Indeed, nilpotent centralizers in $\mathrm{GL}_n$ are connected,
so these simples are indexed by partitions. The ordinary Springer
injection is a bijection because irreducible characteristic-zero
representations of $S_n$ have the same indexing.

For a simple $\mathcal L$, its nonzero inclusion into
$\mathsf I_T^G K$ gives a nonzero map
$\mathsf R_T^G\mathcal L\to K$ by adjunction.
Exactness and finite length show that restriction kills no nonzero
perverse sheaf. It commutes with perverse cohomology, so boundedness
proves the assertion for $\mathcal F$.
Transitivity factors restriction to $T$ through successive proper cuts.
\end{proof}

These statements concern the ordinary nilpotent recipient and its
specified restriction functors. The marked Fourier construction
retains an extension or cotangent variable. Its cut integrals are
different operations, so its positive vanishing theorem survives.

The failure is also visible after a proper forgetting map.
Let $\widetilde{\mathcal N}_2$ consist of pairs $(N,L)$ with
$N(V)\subset L$ and $N(L)=0$. Its projection
$\pi:\widetilde{\mathcal N}_2\to\mathcal N_2$ forgets $L$.
The extension line bundle over $\mathbb P^1$ carries the relative
kernel $\mathcal P$.

\begin{proposition}\label{v4-arith:hall:prop:forget-flag}
The compact cohomology of $f^*R\pi_*\mathcal P$ is nonzero.
Its Frobenius trace is $q(1-q)$, although integration over the
original extension fiber is zero.
\end{proposition}
\begin{proof}
A regular nilpotent has one stable line with nonzero extension
coordinate, giving trace $1$.
At zero, all $q+1$ lines are stable and each contributes $1-q$.
Summing on the marked line therefore gives
$(q-1)+(q+1)(1-q)=q(1-q)\ne0$.
The nonzero trace detects nonzero cohomology.
Including the unipotent factor $K(1)[2]$ divides the trace by $q$
and leaves it nonzero.
\end{proof}

For linear grouping itself, put
$\widetilde{\mathcal P}_V=\mathcal P_V[-1]$.
The inclusion of the product of punctured duals into the punctured
total dual gives
\[
 \widetilde{\mathcal P}_V\boxtimes\widetilde{\mathcal P}_W
 \longrightarrow\widetilde{\mathcal P}_{V\oplus W}.
\]
All parenthesizations include the same locus where every original
covector block is nonzero, so they compose coherently.
The two omitted punctured axes give the triangle
\[
 (\mathcal P_V\boxtimes\mathcal P_W)[-1]
 \longrightarrow\mathcal P_{V\oplus W}
 \longrightarrow
 (\mathcal P_V\boxtimes K_W)\oplus(K_V\boxtimes\mathcal P_W)
 \longrightarrow.
\]
This follows by applying Fourier integration to their localization
triangle, retaining the single shift in
\eqref{v4-arith:hall:eq:fourier-kernel}.

A linear map $g:V\to W$ likewise has the comparison
\[
 \mathsf F_W(Rg_!\mathcal G)
 \simeq(g^\vee)^*\mathsf F_V(\mathcal G).
\]
Both sides integrate the phase $\langle g(v),\eta\rangle$.
Rank normalization adds $[\dim W-\dim V]$ on the right.
These linear comparisons do not construct the nonlinear Hall
incidences or a vanishing-cycle comparison.

\section{Cohomology of the extension correspondence}
\label{v4-arith:hall:sec:coha}

An extension count records the size of a fiber. Equivariant cohomology
also records its characteristic classes. Even the charge-one space
already distinguishes the two constructions: a one-dimensional vector
space has automorphism group $\mathbb C^\times$, and
$H^*(B\mathbb C^\times;\mathbb Q)=\mathbb Q[x]$, with $|x|=2$.
There is one class in every nonnegative even degree.

Let $Q$ be a finite quiver, with vertex set $I$ and $a_{ij}$ arrows
from $i$ to $j$. For a dimension vector $d\in\mathbb N^I$, put
\[
 X_d=\bigoplus_{i,j}\operatorname{Hom}(\mathbb C^{d_i},
                   \mathbb C^{d_j})^{\oplus a_{ij}},\qquad
 G_d=\prod_i GL_{d_i},\qquad \mathfrak M_d=[X_d/G_d].
\]

Fix standard subspaces of dimension $d_i$ in $\mathbb C^{d_i+e_i}$.
Let $X_{d,e}\subset X_{d+e}$ be the representations preserving them,
and let $P_{d,e}\subset G_{d+e}$ preserve the subspaces. The two maps
\[
 \mathfrak M_d\times\mathfrak M_e
 \xleftarrow{\ p\ }[X_{d,e}/P_{d,e}]
 \xrightarrow{\ q\ }\mathfrak M_{d+e}
\]
remember the subrepresentation and quotient, or forget the subspaces.
The map $q$ is proper: its fibers are closed subvarieties of products
of Grassmannians. The map $p$ retains the extension coordinates and
the unipotent automorphisms of the sequence.

Write
\[
 \chi(d,e)=\sum_i d_i e_i-\sum_{i,j}a_{ij}d_i e_j.
\]
The cohomological Hall product acts on
$\bigoplus_d H^*(\mathfrak M_d;\mathbb Q)$. It combines pullback along
$p$, the normal Euler class of $X_{d,e}$, and proper flag pushforward. Its
cohomological degree is $-2\chi(d,e)$.
Since $X_d$ contracts equivariantly to zero,
\[
 H^*(\mathfrak M_d;\mathbb Q)
 =\mathbb Q[x_{i,r}:i\in I,\ 1\leq r\leq d_i]^{\prod_i S_{d_i}}.
\]
Here $x_{i,r}$ are Chern roots of degree two.

\begin{theorem}[Cohomological Hall shuffle product]
\label{v4-arith:hall:thm:coha-shuffle}
For $f$ of charge $d$ and $g$ of charge $e$, the product is
\begin{equation}\label{v4-arith:hall:eq:coha-shuffle}
 f*g=\sum_{\text{vertexwise shuffles}}
 f(x)g(y)
 \frac{\displaystyle\prod_{i,j}\prod_{r=1}^{d_i}
                    \prod_{s=1}^{e_j}(y_{j,s}-x_{i,r})^{a_{ij}}}
      {\displaystyle\prod_i\prod_{r=1}^{d_i}
                    \prod_{s=1}^{e_i}(y_{i,s}-x_{i,r})}.
\end{equation}
The sum is a symmetric polynomial. It defines an associative unital
product. The unit is $1\in H^*(\mathfrak M_0)=\mathbb Q$.
\end{theorem}
\begin{proof}
The normal directions to the invariant-subspace locus are maps from
the subspace at $i$ to the quotient at $j$, one for each arrow
$i\to j$. Their equivariant Euler class is the numerator.
The tangent directions of the flag variety give the denominator.

The fixed points of its maximal torus are precisely the choices of
$d_i$ roots at each vertex. Equivariant localization of the proper
pushforward therefore gives \eqref{v4-arith:hall:eq:coha-shuffle}.
The pushforward is defined before localization, so its sum is polynomial.
This is the extension-correspondence calculation of
Kontsevich--Soibelman, arXiv:1006.2706, \S\S2.2--2.4.

For three factors, either parenthesization chooses the same ordered
three-block partition at every vertex. Every pair of blocks contributes
the same Euler quotient exactly once. The terms, including their
degrees, agree. An empty block contributes no Euler factors and has
one shuffle. This proves associativity and the unit laws.
\end{proof}

Suppose now that $a_{ij}=a_{ji}$. The charge shift
\[
 A_d=H^*(\mathfrak M_d;\mathbb Q)[-\chi(d,d)]
\]
makes multiplication degree zero. Thus a polynomial of degree $r$
has cohomological degree $2r+\chi(d,d)$. Exchanging its two blocks
in \eqref{v4-arith:hall:eq:coha-shuffle} gives
$f*g=(-1)^{\chi(d,e)}g*f$.
We first use the symmetric tensor convention whose crossing on
charges $d,e$ is $(-1)^{\chi(d,e)}$.

There is also an ordinary super convention. Put
\[
 \tau(d,e)=\chi(d,e)+\chi(d,d)\chi(e,e)\pmod2
\]
and choose a bilinear form $\psi$ modulo two with
$\psi(d,e)+\psi(e,d)=\tau(d,e)$.
Such a form exists because $\tau$ is alternating.
The product $f\star g=(-1)^{\psi(d,e)}f*g$ is supercommutative
for the usual cohomological parity. Its associativity follows from
bilinearity of $\psi$. The choice changes signs, not the graded spaces.

\begin{example}\label{v4-arith:hall:ex:zero-one-loop}
For one vertex with no arrows the kernel in
\eqref{v4-arith:hall:eq:coha-shuffle} is $1/(y-x)$.
Let $e_r=x^r$ in charge one. Then
\[
 e_r*e_s=\frac{x^r y^s-x^s y^r}{y-x},\qquad |e_r|=2r+1.
\]

The algebra is the exterior algebra on all $e_r$, $r\geq0$.
Indeed, the product of $e_{r_1},\ldots,e_{r_n}$ is, up to sign,
their alternant divided by the Vandermonde determinant.
For strictly increasing nonnegative exponents these quotients give
the Schur basis of the symmetric polynomials. Repeated exponents
give zero. This proves both spanning and linear independence.

For one vertex with one loop the kernel is $1$ and $|e_r|=2r$.
Products are symmetrized monomials. Their multiplicity factorials
are invertible over $\mathbb Q$, so the algebra is the polynomial
algebra $\mathbb Q[e_0,e_1,\ldots]$.
Neither example has only one generator in charge one.
\end{example}

\section{Primitive fields for a symmetric quiver}
\label{v4-arith:hall:sec:vertex-envelope}

The dual of the Hall product is a coproduct. In the symmetric,
zero-potential case it belongs to a vertex algebra with a precise
primitive description. All duals in this section are graded duals:
\[
 V=\bigoplus_{d,k}(A_d^{-k})^\vee.
\]
Each summand is finite dimensional. The full linear dual, with its
unrestricted product over degrees, is a different recipient.

Use vertex algebras with lower-truncated fields
$Y(a,z)=\sum_n a_{(n)}z^{-n-1}$, vacuum $1$, and translation $D$.
They satisfy $Y(1,z)=\mathrm{id}$,
$Y(Da,z)=[D,Y(a,z)]=\partial_zY(a,z)$ and
$Y(a,z)1=e^{zD}a$, together with mutual locality.
For a pair of states, locality of order $N\geq0$ is equivalent
to $a_{(n)}b=0$ for $n\geq N$.
The symmetric tensor convention is the charge crossing fixed above.

Let $e_i$ denote the $i$th standard dimension vector.
Start with states $v_i$ of charge $e_i$, parity
$\chi(e_i,e_i)\bmod2$ and cohomological degree $-\chi(e_i,e_i)$.
Impose locality orders $N_{ij}=\max\{0,-\chi(e_i,e_j)\}$.
The principal free vertex algebra $P_Q$ is the universal vertex
algebra with these generating states and relations.


The singular products $a_{(n)}b$, $n\geq0$, and translation
form a \emph{vertex Lie algebra}. Its axioms are the sesquilinearity,
skew-symmetry and Jacobi identities of the singular part of a field.
Its universal enveloping vertex algebra $U^{\mathrm{ch}}(P)$
is generated by $P$ with these prescribed singular products and
translation. The coefficient-mode construction appears in the proof.

\begin{theorem}[Dotsenko--Mozgovoy]
\label{v4-arith:hall:thm:vertex-envelope}
For a finite symmetric quiver with zero potential and rational
coefficients, $V$ is canonically a connected cocommutative vertex
bialgebra and $P_Q\simeq V$. Its coproduct is dual to the Hall
product. If
\[
 P=\{a\in V:\Delta a=a\otimes1+1\otimes a\},
\]
then $P$ is a vertex Lie algebra and
$U^{\mathrm{ch}}(P)\simeq V$. Moreover, $D$ has cohomological
degree $-2$, and
\[
 P\simeq W\otimes\mathbb Q[D],\qquad W=P/DP,
 \qquad \dim W_d<\infty.
\]
Every nonzero $W_d^k$ has $k\equiv\chi(d,d)\pmod2$.

After choosing homogeneous lifts of $W$, graded duality gives
\begin{equation}\label{v4-arith:hall:eq:free-coha}
 A\simeq\operatorname{Sym}_{\chi}
                  (W^\vee\otimes\mathbb Q[u]),\qquad |u|=2,
\end{equation}
as algebras. With the product $\star$, this is a free
supercommutative algebra in the ordinary parity convention.
The choice of lifts makes this last presentation noncanonical.
\end{theorem}
\begin{proof}
We recall the maps and the separate finiteness argument in
Dotsenko--Mozgovoy, arXiv:2108.10338, Proposition~5.3 and
Theorems~5.6--5.7. For a sequence $i_1,\ldots,i_m$ of vertices
with total charge $d$, set $N(i,j)=-\chi(e_i,e_j)$. The map
from $P_Q^\vee$ to symmetric polynomials is
\[
 \xi\longmapsto
 \prod_{r<s}(z_r-z_s)^{N(i_r,i_s)}
 \bigl\langle\xi,Y(v_{i_1},z_1)\cdots
                         Y(v_{i_m},z_m)1\bigr\rangle.
\]


Its value is symmetric under simultaneous permutations of vertices
and variables. The vacuum removes negative powers of the last variable. Symmetry
removes them from every variable. The graded functional has finite
degree support, so the expression is polynomial.


When $N(i,i)=-1$, the state is odd. Self-locality of order zero
already implies self-locality of order $-1$.

The map is injective because coefficients of the products of
generating fields span $P_Q$. In the lattice realization, the
creation part of these products is an exponential whose coefficients
are monomials in power sums. These generate the symmetric polynomial
ring over $\mathbb Q$, proving surjectivity.


Declare the $v_i$ primitive. Expand the coproduct by assigning each factor of the ordered field
product to one of two blocks. The cross factor is exactly the Euler
quotient in
\eqref{v4-arith:hall:eq:coha-shuffle}. Thus the comparison preserves
the coproduct, not only graded dimensions.

For primitive $a,b$ of charges $d,e$, compatibility of the fields
and coproduct gives
\[
 \begin{split}
 \Delta(Y(a,z)b)={}&Y(a,z)b\otimes1+1\otimes Y(a,z)b\\
 &+e^{zD}a\otimes b+(-1)^{\chi(d,e)}b\otimes e^{zD}a.
 \end{split}
\]
The two mixed terms are regular. Hence every singular coefficient
$a_{(n)}b$, $n\geq0$, is primitive. Translation preserves primitives
as well.
These operations give the vertex Lie algebra $P$.


Positive charge makes the coalgebra connected: more than $\sum_i d_i$
positive factors annihilate its iterated reduced coproduct in charge $d$.
The connected vertex Milnor--Moore theorem then identifies
$U^{\mathrm{ch}}(P)$ with $V$.

For clarity, the coefficient Lie algebra of $P$ is
\[
 \mathfrak g_P=
 P\otimes\mathbb Q[t,t^{-1}]/(D+\partial_t),\qquad
 [a[m],b[n]]=\sum_{j\geq0}\binom mj(a_{(j)}b)[m+n-j].
\]
Its negative part has
$a[-r]=(D^{r-1}a)[-1]/(r-1)!$ for $r\geq1$.
PBW symmetrization identifies the coalgebra of the vacuum induced
module with $\operatorname{Sym}_{\chi}(\mathfrak g_P^-)$.
In particular, this coalgebra statement concerns every negative mode.

The mode presentation also has an operator $T$, defined on generators
by $T(v_i[n])=-v_i[n+1]$, with
$[D,T]=|d|\,\mathrm{id}$ in positive charge $d$.
Its weights are bounded below, with $D$ raising and $T$ lowering
weight. After dividing $T$ by $|d|$, the following elementary argument
proves that $D$ acts freely. If $Da=0$, the finite span of
$a,Ta,T^2a,\ldots$ is stable under $D,T$. The trace of $[D,T]=1$
would equal its dimension and also vanish, so $a=0$.

Choose a graded complement to $DP$. Induction on weight gives
spanning by its $D$-iterates. Cancellation of the lowest power of
$D$ proves their independence.

This argument alone does not prove $\dim W_d<\infty$.
The framed CoHA character identity in the cited theorem implies
that $(1+q+\cdots+q^{w\cdot d-1})\operatorname{ch}(W_d)$
is a Laurent polynomial, for a framing $w$ with $w\cdot d>0$.
The series $\operatorname{ch}(W_d)$ is bounded below and has
nonnegative coefficients. The product can have finite support only
if the series does. The lattice weights give the asserted parity.
Finally, degreewise duality of the PBW coalgebra, with the chosen
lifts of $W$, proves \eqref{v4-arith:hall:eq:free-coha}.
\end{proof}

The algebra conclusion is also Efimov's freeness theorem,
arXiv:1103.2736, Theorem~1.1. The vertex construction explains
its translation strings. The quotient $P/DP$ has a Lie bracket
induced by the zeroth product. It must not be identified with a
cohomological BPS Lie algebra merely because their graded characters
give the same invariants.

\section{A potential and the perverse filtration}
\label{v4-arith:hall:sec:critical}

A potential changes the coefficients of the correspondence.
Let $W\in\mathbb C Q/[\mathbb C Q,\mathbb C Q]$ be a finite
linear combination of cyclic paths. It defines the regular function
$f_d=\operatorname{Tr}(W)$ on $\mathfrak M_d$.
The Jacobi algebra is the path algebra modulo the cyclic
derivatives $\partial_aW$ for arrows $a$. Its representation
moduli form the critical locus of $f_d$. Critical cohomological Hall theory uses
vanishing cycles of $f_d$, retaining their monodromy.

Choose numbers $\zeta_i$ with positive imaginary parts. For $d\ne0$, set
\[
 Z_\zeta(d)=\sum_i\zeta_i d_i,\qquad
 \mu_\zeta(d)=-\frac{\operatorname{Re}Z_\zeta(d)}{\operatorname{Im}Z_\zeta(d)}.
\]
A representation is $\zeta$-semistable when every nonzero proper
subrepresentation has slope at most that of the full representation.
Fix a slope $\mu$.
Write $\Lambda^\zeta_\mu$ for zero together with the dimension
vectors of that slope, and $\mathfrak M_d^{\zeta\text{-ss}}$
for the stack of semistable representations.

Require $\zeta$ to
be \emph{$\mu$-generic}: $\chi(d,e)=\chi(e,d)$ for all
$d,e\in\Lambda^\zeta_\mu$.
Fix a Serre subcategory $\mathcal S$, closed under subobjects,
quotients and extensions. Require its semistable objects to form
a locally closed substack
$\mathfrak M_d^{\zeta\text{-ss},\mathcal S}$.
These choices enter both the Hall algebra and its primitive term.

We use rational monodromic mixed Hodge modules. This coefficient
category retains the Hodge and weight filtrations together with
vanishing-cycle monodromy. Put
$\mathbb L=H_c(\mathbb A^1;\mathbb Q)$ and use its monodromic
square root $\mathbb L^{1/2}$, of cohomological degree one.

Let $\phi_f^{\mathrm{mon}}$ be the perverse-normalized monodromic
vanishing-cycle functor. Assume that every critical point of $f_d$
lies over zero. Without this assumption, replace it by the sum
of vanishing cycles over all critical values.

Define
\[
 \mathcal A_{W,\mu}^{\zeta,\mathcal S}
 =\bigoplus_{d\in\Lambda^\zeta_\mu}
 H_c\left(\mathfrak M_d^{\zeta\text{-ss},\mathcal S},
 \phi_{f_d}^{\mathrm{mon}}\mathbb Q_{\mathfrak M_d^{\zeta\text{-ss}}}
 \otimes\mathbb L^{-\dim\mathfrak M_d^{\zeta\text{-ss}}/2}
 \right)^\vee.
\]

The functor and dimension shift are taken on the ambient smooth
semistable stack before restriction to $\mathcal S$.
Stack cohomology uses equivariant finite-dimensional approximations.
The dual is the degreewise dual of compactly supported cohomology.
The extension correspondence and Thom--Sebastiani isomorphism
define the critical Hall multiplication on this object.

Let $M_d^{\zeta\text{-ss}}$ be the coarse moduli space.
Its normalized intersection complex $\operatorname{IC}_d$ is
Verdier self-dual. On a smooth component it is
$\mathbb L^{-\dim M_d/2}\mathbb Q_{M_d}$.
Define the BPS sheaf to be
$\phi_{f_d}^{\mathrm{mon}}\operatorname{IC}_d$ if the stable
locus is nonempty, and zero otherwise.

Denote by
$\mathcal B_{W,\mu}^{\zeta,\mathcal S}$ the direct sum of the
duals of its compactly supported cohomology restricted to the
coarse loci selected by $\mathcal S$.
The remaining universal factor is
\[
 H(B\mathbb C^\times)_{\mathrm{vir}}
 =\mathbb L^{1/2}\otimes H^*(B\mathbb C^\times;\mathbb Q).
\]
Its half-Tate shift is part of the normalization.

\begin{theorem}[Davison--Meinhardt]
\label{v4-arith:hall:thm:critical-pbw}
With these hypotheses there is a PBW isomorphism of
charge-graded, cohomologically graded monodromic mixed Hodge objects
\[
 \operatorname{Sym}_{\tau}\left(
 H(B\mathbb C^\times)_{\mathrm{vir}}\otimes
 \mathcal B_{W,\mu}^{\zeta,\mathcal S}\right)
 \xrightarrow{\ \sim\ }\mathcal A_{W,\mu}^{\zeta,\mathcal S}.
\]
Here the symmetry is the Koszul sign multiplied by
$(-1)^{\tau(d,e)}$, with $\tau$ as above on
$\Lambda^\zeta_\mu$.
The perverse associated graded of the Hall algebra is free
commutative for this symmetry. After a bilinear sign modification
$\psi+\psi^{\mathrm T}=\tau$ of multiplication, that associated
graded is free supercommutative for the ordinary Koszul symmetry.
The displayed PBW map concerns underlying graded objects. It need
not be an algebra isomorphism for the original Hall multiplication.
\end{theorem}

This is the absolute form of Theorems~A and~C of
Davison--Meinhardt, arXiv:1601.02479. The perverse filtration comes
from the map from the semistable stack to its coarse moduli space.
It is not the filtration by total dimension alone. The PBW map uses
the BPS embedding and symmetrized Hall multiplication. The original
algebra can therefore have nonzero commutators whose leading
perverse symbols commute.

The recipient allows infinitely many cohomological degrees.
For each fixed charge and weight, only finitely many cohomological
degrees occur, and the weights are bounded below at fixed charge.
Symmetric powers are formed by sums in this recipient. They are
not unrestricted products of all cohomological degrees.
The zero-potential algebra theorem
\eqref{v4-arith:hall:eq:free-coha} retains its stronger conclusion
in its own finite symmetric-quiver setting.

\section{An ordinary coproduct obstruction on a shuffle algebra}
\label{v4-arith:hall:sec:cooperation}

The preceding vertex coproduct is dual to a symmetric Hall product.
A deformed shuffle multiplication poses a different compatibility
problem. Its charge-one elements can already prevent an ordinary
counital coproduct from being multiplicative.

Work over
\[
 R=\mathbb Q[h_1,h_2],\qquad h_3=-h_1-h_2,\qquad
 a=-(h_1^2+h_1h_2+h_2^2),\qquad b=-h_1h_2(h_1+h_2).
\]
Put $N(v)=\prod_i(v+h_i)=v^3+av+b$. For $j=0,1$, let
\[
 \mathcal A^{(j)}=\bigoplus_{n\geq0}R[x_1,\ldots,x_n]^{S_n},
 \qquad K_j(v)=v^{2j-1}N(v).
\]

These are algebraic direct sums. Charge is $n$. All cochain degrees
and parities are zero. There is no completion or inversion of
parameters. The product is
\begin{equation}\label{v4-arith:hall:eq:deformed-shuffle}
 (F*_jG)(X)=\sum_{I\sqcup J=[d+e],\ |I|=d}
 F(X_I)G(X_J)\prod_{i\in I,k\in J}K_j(x_k-x_i).
\end{equation}

There is no factorial denominator. For $j=0$, opposite residues
cancel along each true diagonal, so the sum is polynomial.
For $j=1$ every summand is polynomial. Ordered three-block
partitions prove associativity in either case. The unit is
$1_0\in A_0=R$, distinct from the constant polynomial in charge one.

The binary image remembers the coefficient ring. Write
$c=(x+y)/2$, $v=y-x$, $t=v^2$ and $B=R[c,t]$. Then
$R[x,y]=B\oplus vB$, and the multiplication map is
\begin{equation}\label{v4-arith:hall:eq:binary-image}
 \mu_j(A+vC)=2t^j\bigl((t+a)A+bC\bigr).
\end{equation}

Consequently its image is $t^j(t+a,b)$. The complex
$[B^2\xrightarrow{\mu_j}B]$ in degrees $-1,0$ has
\[
 H^{-1}=B\,(b,-(t+a)),\qquad
 H^0=B/t^j(t+a,b).
\]

To compute the kernel, reduce $(t+a)A+bC=0$ modulo $t+a$.
The quotient $R[c]$ is a domain, so $C$ is divisible by $t+a$.
Cancellation gives the displayed generator. Multiplication by
$t^j$ adds no kernel.


The two cokernels have different parameter
annihilators:
\[
 \operatorname{Ann}_R B/(t+a,b)=(b),\qquad
 \operatorname{Ann}_R B/t(t+a,b)=0.
\]
For the first, set $t=-a$. For the second, evaluation at $t=0$
maps the quotient onto the faithful $R$-module $R[c]$.
Thus the two binary lattices cannot be interchanged over $R$
by an operation-preserving charge-graded equivalence.

\begin{theorem}\label{v4-arith:hall:thm:no-ordinary-coproduct}
Neither $\mathcal A^{(j)}$ admits an $R$-linear, charge-preserving,
unital algebra map
\[
 \Delta:\mathcal A^{(j)}\longrightarrow
          \mathcal A^{(j)}\otimes_R^{\mathbb L}\mathcal A^{(j)}
\]
with ordinary Koszul multiplication and two-sided counit given
by projection to charge zero. The same holds for the subalgebra
generated by charge one. No coassociativity hypothesis is needed.
\end{theorem}
\begin{proof}
Set $e_m=x^m\in A_1$. The difference of the two kernel orders gives
\[
 [f,g]_{*_j}=2bt^j\frac{f(x)g(y)-g(x)f(y)}{y-x}.
\]
In particular,
\[
 \begin{aligned}
 [e_0,e_3]&=2bt^j(x^2+xy+y^2),&
 [e_1,e_2]&=2bt^jxy,\\
 [e_0,e_1]&=2bt^j,&e_0*_je_0&=2t^j(t+a).
 \end{aligned}
\]
They satisfy the exact relation
\begin{equation}\label{v4-arith:hall:eq:four-current-relation}
 [e_0,e_3]-3[e_1,e_2]+a[e_0,e_1]=b(e_0*_je_0).
\end{equation}

Charge preservation and both counits force
$\Delta e_m=e_m\otimes1_0+1_0\otimes e_m$.
In the ordinary tensor product, mixed terms in the commutator
of two even primitive elements cancel. The $(1,1)$ component of
the coproduct of the left side of
\eqref{v4-arith:hall:eq:four-current-relation} is therefore zero.
The right side gives $2b\,e_0\otimes e_0$, which is nonzero
in $R[x,y]$.

Every ambient charge component is free over $R$, so its derived
tensor is ordinary. For the charge-one generated subalgebra the
same detecting component is still the free tensor $R[x,y]$.
Taking cohomology preserves the contradiction. At
$(h_1,h_2,h_3)=(1,1,-2)$ its coefficient is $2b=-4$.
The proof neither divides by $b$ nor asserts persistence at $b=0$.
\end{proof}

The tensor in this theorem is essential. A rational replacement
can carry a compatible cooperation. For ordered blocks $X,Y$, set
\[
 Q_{d,e}^{(j)}(X,Y)=\prod_{x\in X,y\in Y}K_j(y-x),\qquad
 \Delta_{d,e}^{\mathrm{rat}}F=
                \frac{F(X,Y)}{Q_{d,e}^{(j)}(X,Y)},\qquad
 r(v)=\frac{v^3+av-b}{N(v)}.
\]

Use polynomial rings localized at all spectral differences and
their translates $v\pm h_i$, then take the blockwise symmetric
invariants. No parameter is inverted by itself, but closed spectral
loci have been removed.
Take $F$ of bidegree $(d,e)$ and $G$ of bidegree $(d',e')$.
Their target product is the sum of
\[
 \begin{split}
 &F(X_1,Y_1)G(X_2,Y_2)
   \prod_{x\in X_1,x'\in X_2}K_j(x'-x)\\
 &\hspace{8mm}\cdot\prod_{y\in Y_1,y'\in Y_2}K_j(y'-y)
   \prod_{y\in Y_1,x'\in X_2}r(y-x')
 \end{split}
\]
over the two prescribed ordered splittings of $X$ and $Y$.
The last factor changes the interchange of the middle two blocks.

\begin{proposition}\label{v4-arith:hall:prop:rational-coproduct}
This target multiplication is associative and unital.
The maps $\Delta^{\mathrm{rat}}$ are multiplicative,
coassociative in the corresponding three-block localizations,
and counital at charge zero.
\end{proposition}
\begin{proof}
In a shuffle term, division by the total cross kernel cancels
the factors between $X_1$ and $Y_2$. The remaining middle ratio is
$K_j(x'-y)/K_j(y-x')=r(y-x')$, for $y\in Y_1$, $x'\in X_2$.
This is exactly the displayed target product.

For three factors, either parenthesization contains the same
internal kernels and the same earlier-right, later-left exchange
factors. This proves associativity. Iterating cooperation divides
by $Q_{12}Q_{13}Q_{23}$ in both orders. Empty products give the
unit and both counits.
\end{proof}

For \eqref{v4-arith:hall:eq:four-current-relation}, the new mixed
term satisfies
\[
 (1-r(v))v(v^2+a)=b(1+r(v)),
\]
as follows directly from $1-r(v)=2b/N(v)$.
It supplies the term absent from the ordinary tensor calculation.


Nevertheless,
\[
 \Delta^{\mathrm{rat}}_{1,1}(1_2)=
 \begin{cases}v/N(v),&j=0,\\1/(vN(v)),&j=1,
 \end{cases}
\]
is not polynomial. Localizing an equation that vanishes on a
shifted collision annihilates that closed restriction.
Extending this rational cooperation to supported coefficients remains
a separate problem. Such an extension must retain the integral
images \eqref{v4-arith:hall:eq:binary-image} and derived collision complexes.
An ordinary coefficient bialgebra and a strong coalgebra
for cut convolution have different tensors. The theorem excludes
the latter only when its marked comparison induces the former.

\section{The arithmetic Hall comparison}
\label{v4-arith:hall:sec:arithmetic-comparison}

The constructions above specify several recipients of Hall data.
The scalar Hall algebra retains automorphism-weighted extension
counts. The Fourier kernel retains an attaching morphism and
compatible primitive contractions. Cohomological Hall theory retains
characteristic classes. A potential additionally retains vanishing
cycles and their monodromy. A comparison must name which of these
operations it preserves.

\begin{question}\label{v4-arith:hall:qn:arithmetic}
Construct an independently defined arithmetic Hall category and
an operation-preserving functor to the intended geometric or field
category. Its extension correspondences must induce the specified
Hall multiplication and primitive cut maps. It must retain automorphism and Tate normalizations, place labels,
substitutions and nonconsecutive collision supports. The empty unit
and marked local vacuum require their own compatible maps.
All comparisons must respect composition.


For a critical CoHA target, construct the arithmetic quiver,
potential, stability data and normalized coefficient comparison.
For a field target, construct the actual fields, translation and
supported products on the proposed states.
\end{question}

The flag kernels solve the primitive cut problem with the flag
retained. Their descent failures locate the additional datum needed
by an unmarked recipient. The symmetric-quiver theorem constructs
primitive fields on its graded dual. Neither statement constructs
an arithmetic Hall functor or an equivalence with an independently
specified arithmetic category.

The formal Euler expansion also keeps a precise scope. On
$\operatorname{Re}s>1$,
\[
 \log\zeta(s)=\sum_p\sum_{k\geq1}\frac{p^{-ks}}k,\qquad
 -\frac{\zeta'(s)}{\zeta(s)}
       =\sum_p\sum_{k\geq1}(\log p)p^{-ks}.
\]

The coefficients of the first series are rational. Logarithms
appear after differentiation. Neither series specifies a finite
cyclic word in a path algebra. A potential requires such a cyclic
class, or a specified completion and a trace map on it. A class
in cyclic homology or an exponential Hodge recipient would require
its own construction and comparison with the Euler series.
The scalar identity supplies no canonical lift of this kind.
